\clearpage
\section{Introduction}
\label{sec:Intro}

Higher education is widely regarded as one of the most effective instruments for promoting intergenerational mobility and reducing earnings inequality. Yet the extent to which expanding access to universities translates into equitable labor market gains across demographic groups and geographic regions remains an open empirical question, particularly in developing economies characterized by deep structural dualism. Brazil---famously described as \emph{Belíndia}, a country that simultaneously resembles Belgium and India \cite{bacha1976}---offers a compelling setting in which to study this question: its labor market exhibits pronounced wage gaps along racial, gender, and regional lines that persist even after controlling for observable human capital \cite{AriasYamadaTejerina2002, Telles2004, deoliveira2021}.

In 2005, the Brazilian federal government launched the \acrfull{prouni}, a large-scale scholarship program that leverages tax exemptions for private \acrfull{hei} to provide tuition-free university seats to low-income students. By 2019, ProUni had allocated over 3.5 million scholarships, with built-in affirmative action mandates reserving a share for self-declared non-white applicants proportional to their state-level population share \cite{MEC2024}. The program thus constitutes both a human capital supply shock and a race-conscious redistributive policy, making it a natural experiment for studying how educational expansion reshapes labor market inequality.

Despite its scale and policy salience, the causal labor market effects of ProUni remain insufficiently understood. While the existing literature has documented positive effects of affirmative action in elite public universities \cite{RiehlMachadoReyes2023} and estimated aggregate returns to education in Brazil using Mincerian frameworks \cite{Card1999, Mincer1974}, no study to date has exploited the staggered, spatially heterogeneous rollout of ProUni to identify its causal impact on formal sector wages at the local labor market level, let alone its distributional effects across gender, race, and region.

This thesis addresses this gap by asking: \emph{How does the expansion of ProUni scholarships affect formal sector earnings, and do these effects vary systematically by gender, ethnicity, and geographic region?}

To answer this question, I construct a novel panel dataset linking three administrative data sources---ProUni scholarship records, the \acrfull{rais} employer--employee census, and Census population data---at the microregion level over the period 2005--2019. The key innovation in the research design is the construction of a \emph{lagged cumulative scholarship density} measure that captures the local ``dose'' of ProUni-educated graduates entering each microregion's labor market, normalized by the youth population. This continuous treatment variable is discretized into dose strata to test for non-linear effects, specifically the hypothesis that low-dose regions experience ``scarcity rents'' (positive wage premia from a scarce supply of skilled workers) while high-dose regions experience ``saturation'' effects (credential inflation from oversupply).

The identification strategy employs the heterogeneity-robust \acrfull{did} estimator of \citeA{CallawaySantAnna2021}, which avoids the well-documented biases of conventional \acrfull{twfe} regression in settings with staggered treatment adoption and heterogeneous effects \cite{GoodmanBacon2021, deChaisemartin2020}. This approach estimates group--time average treatment effects for each treatment cohort and aggregates them into dynamic event studies, allowing for transparent inspection of pre-trends and the evolution of treatment effects over time. The continuous treatment extension follows \citeA{CallawayGoodmanBaconSantAnna2024}.

The thesis makes three contributions to the literature. First, it provides the first causal estimates of ProUni's labor market effects using a design that exploits spatial and temporal variation in program intensity, moving beyond the descriptive and selection-on-observables approaches that have characterized prior work. Second, by estimating separate models for non-white and white workers, men and women, and structurally distinct regions (North/Northeast versus South/Southeast), it reveals the distributional anatomy of a universal education policy---documenting where and for whom returns materialize and where structural barriers attenuate them. Third, the dose-stratified design contributes to the growing methodological literature on continuous treatments in difference-in-differences frameworks by providing an applied demonstration of how discretization choices interact with substantive economic hypotheses about general equilibrium effects.

The remainder of the thesis is organized as follows. Section~\ref{sec:literature} reviews the theoretical foundations in human capital and segmented labor market theory, as well as the empirical evidence on returns to education in Brazil and internationally. Section~\ref{sec:institutional_context} provides the institutional context, covering Brazil's higher education system, the design and legal foundations of ProUni, and concurrent policy interventions. Section~\ref{sec:methodology} details the data sources, treatment construction, identification assumptions, and estimation strategy. Section~\ref{sec:results} presents the main results. Section~\ref{sec:discussion} discusses mechanisms and policy implications. Section~\ref{sec:conclusion} concludes.
